\scanpage{159}
\uline{Formes définies positives sur les motifs}

Ceci implique déjà que les $\operatorname{End}(h^i(X))$ sont des \uline{algèbres de Riemann}, et par suite $\frac12$ simples, (ce qui permet de construire la catégorie des des motifs comme catégorie \uline{abélienne}).

De plus, cette catégorie des motifs $\mathcal M$, et sa complétée $\widetilde{\mathcal M}$, \unsure{recevra} d'importantes structures suppl.

Discutons d'abord le degré d'indétermination possible du système $\varphi_X$. Pour ceci, \unsure{notons} \deleted{pour} regardons ce qui se passe si on remplace les $(\varphi_i)$ par $(\varepsilon_i\varphi_i)$ ($\varepsilon_i$, \uline{signes} ne dépendant que de $i$). Les axiomes (i) et (iii) restent vérifiés, mais (ii) seulement si on a $\varepsilon_i\varepsilon_j=\varepsilon_{\unsure{i+j|ij}}$, donc $\varepsilon_i$ doit \deleted{être de la forme $\unsure{(\pm1)}$} $+1$ pour tt $i$, \uline{ou} $\unsure{(-1)^i|\varepsilon^i}$ --- il est caractérisé en se donnant ce qu'il doit être en degré $1$.
On imposera, dans le cas de la géom. alg., qu'il corresponde aux formes de polarisation sur les V.A.
\mspage{158}
